\documentclass[11pt,a4paper]{article}

% ---------------------------------------------------------------------------
% Paquetes
% ---------------------------------------------------------------------------
\usepackage[utf8]{inputenc}
\usepackage[T1]{fontenc}
\usepackage{lmodern}
\usepackage[spanish]{babel}
\usepackage[margin=2.5cm]{geometry}
\usepackage{amsmath}
\usepackage{graphicx}
\usepackage{float}
\usepackage{booktabs}
\usepackage{array}
\usepackage{xcolor}
\usepackage{hyperref}

\hypersetup{colorlinks=true, linkcolor=blue!50!black, urlcolor=blue!50!black}

% ---------------------------------------------------------------------------
% Atajos
% ---------------------------------------------------------------------------
\newcommand{\dato}[1][2cm]{\underline{\hspace{#1}}}
\newcommand{\us}{\ensuremath{\mu}s}

% Placeholder de captura: el documento COMPILA sin las imagenes. Cuando
% tengas la foto del osciloscopio, reemplazar el recuadro por el
% \includegraphics que figura escrito adentro.
\newcommand{\captura}[2]{%
  \begin{figure}[H]
  \centering
  \fbox{\parbox[c][4.5cm][c]{0.82\textwidth}{\centering
    \itshape Captura pendiente\\[6pt]
    \normalfont\small #2\\[10pt]
    \footnotesize Reemplazar este recuadro por:\\
    \texttt{\textbackslash includegraphics[width=0.82\textwidth]\{img/#1.png\}}}}
  \caption{#2}
  \label{fig:#1}
  \end{figure}}

\title{\textbf{Medición de Tiempos de Interrupción}\\[4pt]
\large Sistema de Control de Acceso --- FRDM-K64F}
\author{Santiago Riverós \and Ignacio Sammartino \and Rodrigo Devesa \and Leandro Yiu}
\date{TP1 --- Sistemas Embebidos}

\begin{document}

\maketitle

\begin{abstract}
\noindent
Este informe caracteriza experimentalmente los tiempos de interrupción del
firmware mediante un pin de instrumentación medido con osciloscopio. El
objetivo es verificar que el sistema cumple sus restricciones de tiempo
real, cuantificar el costo de las abstracciones de la arquitectura en capas,
y validar la asignación de prioridades de interrupción entre el SysTick y
el lector de tarjetas.
\end{abstract}

\tableofcontents
\newpage

% =============================================================================
\section{Método}
% =============================================================================

El firmware incluye un módulo de instrumentación (\texttt{test\_isr\_cfg.h})
que pone un pin dedicado en alto al entrar a una rutina y en bajo al salir.
En el osciloscopio, el \textbf{ancho del pulso} es el tiempo de ejecución de
la rutina y el \textbf{período entre flancos de subida} es cada cuánto se
ejecuta. La instrumentación se resuelve en tiempo de compilación: sólo la
rutina seleccionada genera código, así que medir una rutina no afecta el
tiempo de las demás.

\begin{table}[H]
\centering
\begin{tabular}{@{}ll@{}}
\toprule
\textbf{Parámetro} & \textbf{Valor} \\
\midrule
Microcontrolador        & MK64FN1M0VLL12 (Cortex-M4) \\
Frecuencia de core      & 100 MHz (1 ciclo = 10 ns) \\
Optimización            & \texttt{-O0} (build de Debug) \\
Frecuencia del SysTick  & 1 kHz (período 1 ms) \\
Pin de medición         & PTC17 \\
Osciloscopio            & \dato[4cm] \\
\bottomrule
\end{tabular}
\caption{Condiciones bajo las que se realizaron todas las mediciones.}
\label{tab:condiciones}
\end{table}

Todos los tiempos corresponden a una compilación sin optimizar. Con
\texttt{-O2} los valores absolutos serían menores; las relaciones entre
ellos, que sostienen las conclusiones de este informe, se mantienen.

\newpage
% =============================================================================
\section{Mediciones realizadas}
% =============================================================================

Se realizaron nueve mediciones sobre las distintas fuentes de interrupción y
tareas periódicas del sistema. Cada una activa una única rutina de
instrumentación mediante la constante \texttt{TEST\_ISR\_TARGET}.

% -----------------------------------------------------------------------------
\subsection{M1 --- SysTick completo}
% -----------------------------------------------------------------------------
\textbf{Configuración:} \texttt{TEST\_ISR\_TARGET = TEST\_ISR\_SYSTICK\_TOTAL}

Mide la ejecución completa de \texttt{SysTick\_Handler()}, la interrupción
periódica de mayor jerarquía del sistema: despacha las tres tareas
registradas (refresco del display, actualización del encoder y base de
tiempo en milisegundos). Tiene que completarse siempre antes de que llegue
el próximo tick, así que es la referencia contra la que se comparan las
demás mediciones periódicas.

% -----------------------------------------------------------------------------
\subsection{M2 --- Refresco del display}
% -----------------------------------------------------------------------------
\textbf{Configuración:} \texttt{TEST\_ISR\_TARGET = TEST\_ISR\_DISPLAY\_REFRESH}

Mide \texttt{HAL\_Refresh\_Task()}, la tarea que se espera dominante dentro
del SysTick: hace \textit{bit-banging} de dos tramas de 16 bits (una de
\textit{blanking} anti-\textit{ghosting} y una útil) hacia los
shift-registers 74HC595. Se mide con el display a brillo máximo y a brillo
mínimo, ya que el driver está escrito sin ramas por modo de operación y se
espera que el ancho no cambie entre ambas condiciones.

% -----------------------------------------------------------------------------
\subsection{M3 --- FSM del encoder}
% -----------------------------------------------------------------------------
\textbf{Configuración:} \texttt{TEST\_ISR\_TARGET = TEST\_ISR\_ENCODER\_UPDATE}

Mide \texttt{actualizarEncoder()}, que lee tres pines y avanza una máquina
de estados de antirrebote. El encoder se sondea por \textit{polling} desde
el mismo tick de 1 ms en lugar de usar interrupciones de pin propias; esta
medición permite evaluar el costo real de esa decisión.

% -----------------------------------------------------------------------------
\subsection{M4 --- Base de tiempo en milisegundos}
% -----------------------------------------------------------------------------
\textbf{Configuración:} \texttt{TEST\_ISR\_TARGET = TEST\_ISR\_TIMER\_TICK}

Mide \texttt{Timer\_Tick\_Callback()}, la rutina más simple del sistema (un
incremento de contador). Sirve como referencia del costo mínimo de un
callback registrado en el \textit{scheduler}, incluyendo la llamada
indirecta a través del puntero a función.

% -----------------------------------------------------------------------------
\subsection{M5 --- Costo del despachador del SysTick}
% -----------------------------------------------------------------------------
No requiere una medición directa ni una captura propia: se obtiene
combinando los resultados de M1 a M4. \texttt{SysTick\_Handler()} no ejecuta
las tres tareas en línea, sino que recorre un arreglo de punteros a función
y las invoca; ese bucle de despacho es el precio de que los módulos puedan
registrarse en el \textit{scheduler} sin que el MCAL los conozca. Su costo
se calcula en la Sección~\ref{sec:analisis}.

% -----------------------------------------------------------------------------
\subsection{M6 --- Lector de tarjetas: captura de un bit}
% -----------------------------------------------------------------------------
\textbf{Configuración:} \texttt{TEST\_ISR\_TARGET = TEST\_ISR\_CARD\_CLK}, y
luego \texttt{TEST\_ISR\_PORT\_DISPATCH}.

Mide la ISR más crítica del sistema: \texttt{IRQ\_CardClk()} se ejecuta una
vez por cada bit leído de la banda magnética, y perder una sola ejecución
corrompe la trama completa. Midiendo primero sólo el cuerpo de la ISR
(\texttt{CARD\_CLK}) y luego la interrupción completa con el despachador de
puerto incluido (\texttt{PORT\_DISPATCH}), se aísla el costo del mecanismo
genérico de \texttt{gpioIRQDispatch()}, que recorre los 32 pines del puerto
buscando cuál generó la interrupción.

% -----------------------------------------------------------------------------
\subsection{M7 --- Lector de tarjetas: apertura y cierre del \textit{swipe}}
% -----------------------------------------------------------------------------
\textbf{Configuración:} \texttt{TEST\_ISR\_TARGET = TEST\_ISR\_CARD\_ENABLE}

Mide \texttt{IRQ\_CardEnable()}, que se ejecuta dos veces por pasada de
tarjeta (flanco de bajada al empezar, de subida al terminar). A diferencia
de \texttt{IRQ\_CardClk()}, esta ISR llama a \texttt{gpioIRQ()} para armar y
desarmar la interrupción de CLOCK, lo que implica configurar registros del
periférico (PCR, ISFR, NVIC); se espera un tiempo de ejecución notablemente
mayor.

% -----------------------------------------------------------------------------
\subsection{M8 --- Latencia de interrupción}
% -----------------------------------------------------------------------------
\textbf{Configuración:} \texttt{TEST\_ISR\_TARGET = TEST\_ISR\_CARD\_CLK},
dos canales: CH1 en PTB10 (CLOCK del lector, la señal física) y CH2 en
PTC17 (la respuesta del software).

Mide el tiempo entre el flanco de bajada de CLOCK y el instante en que el
CPU efectivamente entra a la ISR. Esta medición permite validar
experimentalmente la asignación de prioridades del NVIC: en
\texttt{App\_Init()} se le baja la prioridad al SysTick
(\texttt{APP\_NVIC\_PRIO\_SYSTICK = 1}) para que la interrupción de PORTB
(prioridad 0, el lector) siempre pueda interrumpirlo. Si esa política
funciona, la latencia medida debería ser pequeña y no acercarse nunca a la
duración de la ISR del SysTick (M1).

% -----------------------------------------------------------------------------
\subsection{M9 --- Super-loop}
% -----------------------------------------------------------------------------
\textbf{Configuración:} \texttt{TEST\_ISR\_TARGET = TEST\_ISR\_APP\_LOOP}

Mide el período de \texttt{App\_Run()}. No es una interrupción, pero
determina la capacidad de respuesta de la aplicación: \texttt{encoder.h}
documenta que un evento del encoder vive como máximo 1 ms antes de
limpiarse, por lo que el super-loop debe completar una vuelta en menos de
ese tiempo o se pierden pulsaciones. Se mide en reposo y durante la
decodificación de una tarjeta, que es la carga adicional más pesada que
puede caerle al lazo principal.

\newpage
% =============================================================================
\section{Resultados}
% =============================================================================

\subsection{Anchos de pulso y períodos}

\begin{table}[H]
\centering
\begin{tabular}{@{}llccc@{}}
\toprule
\textbf{Ensayo} & \textbf{Rutina / condición} & \textbf{Ancho} & \textbf{Período} & \textbf{Ciclos} \\
\midrule
M1 & \texttt{SysTick\_Handler}                    & \dato[1.4cm]\,\us & \dato[1.3cm]\,ms  & \dato[1.3cm] \\
M2 & \texttt{HAL\_Refresh\_Task} (brillo máx.)     & \dato[1.4cm]\,\us & 1\,ms             & \dato[1.3cm] \\
M2 & \texttt{HAL\_Refresh\_Task} (brillo mín.)     & \dato[1.4cm]\,\us & 1\,ms             & \dato[1.3cm] \\
M3 & \texttt{actualizarEncoder} (reposo)           & \dato[1.4cm]\,\us & 1\,ms             & \dato[1.3cm] \\
M3 & \texttt{actualizarEncoder} (girando)          & \dato[1.4cm]\,\us & 1\,ms             & \dato[1.3cm] \\
M4 & \texttt{Timer\_Tick\_Callback}                & \dato[1.4cm]\,\us & 1\,ms             & \dato[1.3cm] \\
M6 & \texttt{IRQ\_CardClk} (solo la ISR)           & \dato[1.4cm]\,\us & \dato[1.3cm]\,\us & \dato[1.3cm] \\
M6 & \texttt{gpioIRQDispatch} (con despacho)       & \dato[1.4cm]\,\us & \dato[1.3cm]\,\us & \dato[1.3cm] \\
M7 & \texttt{IRQ\_CardEnable} (apertura)           & \dato[1.4cm]\,\us & --- (2 por pasada) & \dato[1.3cm] \\
M7 & \texttt{IRQ\_CardEnable} (cierre)             & \dato[1.4cm]\,\us & --- (2 por pasada) & \dato[1.3cm] \\
M9 & \texttt{App\_Run} (reposo)                    & \dato[1.4cm]\,\us & \dato[1.3cm]\,\us & \dato[1.3cm] \\
M9 & \texttt{App\_Run} (decodificando tarjeta)     & \dato[1.4cm]\,\us & \dato[1.3cm]\,\us & \dato[1.3cm] \\
\bottomrule
\end{tabular}
\caption{Resultados de todas las mediciones de ancho de pulso y período.}
\label{tab:resultados}
\end{table}

\subsection{Latencia de interrupción}

La medición M8 es de naturaleza distinta a las anteriores: no es un ancho de
pulso sino un retardo entre dos canales, así que se reporta aparte.

\begin{table}[H]
\centering
\begin{tabular}{@{}lcc@{}}
\toprule
\textbf{Magnitud} & \textbf{Valor} & \textbf{Ciclos} \\
\midrule
Latencia mínima observada (CH1$\to$CH2) & \dato[1.6cm]\,\us & \dato[1.5cm] \\
Latencia máxima observada (CH1$\to$CH2) & \dato[1.6cm]\,\us & \dato[1.5cm] \\
\bottomrule
\end{tabular}
\caption{M8: latencia de interrupción del lector de tarjetas, sobre una serie de capturas.}
\label{tab:latencia}
\end{table}

\newpage
% =============================================================================
\section{Análisis}
\label{sec:analisis}
% =============================================================================

\subsection{Cumplimiento de la restricción de tiempo real (M1)}

Según la Tabla~\ref{tab:resultados}, el período de \texttt{SysTick\_Handler}
fue de \dato\,ms, lo que confirma que el SysTick está configurado
correctamente a 1 kHz. La ISR consume \dato\,\us\ de cada milisegundo, es
decir un \dato\,\% de la carga de CPU en ese contexto, dejando un
\dato\,\% de margen antes de comprometer el próximo tick. Ese margen
equivale a \dato\,\us: la ISR podría crecer esa cantidad de tiempo antes de
que el sistema empiece a perder ticks.

\subsection{Componente dominante del SysTick (M2)}

El refresco del display consume \dato\,\us, un \dato\,\% del tiempo total
de la ISR medido en M1: es, en efecto, el componente dominante. La
diferencia entre brillo máximo y mínimo fue de \dato\,\us. Dado que el
atenuador está implementado por reparto de turnos de multiplexado --- en
los turnos ``vacíos'' se manda la misma trama con los segmentos apagados,
en lugar de ejecutar menos código ---, este resultado \dato{} (confirma /
no confirma) que el costo de CPU del refresco es independiente del brillo
elegido por el usuario.

\subsection{Costo del \textit{polling} del encoder (M3)}

La FSM del encoder consume \dato\,\us\ por tick, un \dato\,\% del
presupuesto de la ISR del SysTick. Es un costo \dato{} (despreciable /
significativo) frente al del display, lo que respalda sondear el encoder
por \textit{polling} en el tick de 1 ms en lugar de dedicarle una
interrupción de pin propia, que competiría por el NVIC con el lector de
tarjetas.

\subsection{Piso de costo de un callback (M4)}

El contador de milisegundos, la rutina más simple del sistema, consume
\dato\,\us. Este valor representa aproximadamente el costo mínimo de
cualquier callback registrado en el \textit{scheduler}, incluyendo la
llamada indirecta a través del puntero a función.

\subsection{Costo del despachador del SysTick (M5)}

Restando a la duración total medida en M1 la suma de las tres tareas
medidas en M2, M3 y M4 se obtiene el tiempo que insume exclusivamente el
bucle despachador de \texttt{SysTick\_Handler()}:

\[
t_{\text{despachador}} \;=\; t_{\text{M1}} \;-\; \left( t_{\text{M2}} + t_{\text{M3}} + t_{\text{M4}} \right)
\]

\begin{table}[H]
\centering
\begin{tabular}{@{}lcc@{}}
\toprule
\textbf{Componente} & \textbf{Tiempo} & \textbf{\% del total (M1)} \\
\midrule
Refresco del display (M2, brillo máx.) & \dato[1.5cm]\,\us & \dato[1.3cm]\,\% \\
FSM del encoder (M3, reposo)           & \dato[1.5cm]\,\us & \dato[1.3cm]\,\% \\
Base de tiempo (M4)                    & \dato[1.5cm]\,\us & \dato[1.3cm]\,\% \\
\midrule
Suma de las tres tareas                & \dato[1.5cm]\,\us & \dato[1.3cm]\,\% \\
Total medido (M1)                      & \dato[1.5cm]\,\us & 100\,\% \\
\midrule
\textbf{Despachador (diferencia)}      & \dato[1.5cm]\,\us & \dato[1.3cm]\,\% \\
\bottomrule
\end{tabular}
\caption{Descomposición del tiempo de la ISR del SysTick.}
\label{tab:despachador-systick}
\end{table}

El bucle despachador agrega \dato\,\us\ de sobrecarga sobre las tareas, un
\dato\,\% del total. Este es el costo concreto de permitir que
\texttt{App\_Init()} registre tareas con \texttt{HAL\_Scheduler\_AddTask()}
sin que el MCAL conozca su existencia. El resultado \dato{} (justifica / no
justifica) mantener el mecanismo genérico frente a llamar a las tres
funciones en forma directa dentro de la ISR.

\subsection{Costo del despachador de puerto (M6)}

Comparando las dos filas de M6 en la Tabla~\ref{tab:resultados}, el cuerpo
de \texttt{IRQ\_CardClk()} por sí solo consume \dato\,\us, mientras que la
interrupción completa --- con \texttt{gpioIRQDispatch()} incluido ---
consume \dato\,\us. La diferencia,

\[
\Delta_{\text{PORT}} \;=\; \dato\,\us,
\]

es el costo de recorrer las 32 posiciones del puerto para encontrar cuál
generó la interrupción, aunque sólo una tenga el \textit{flag} activo.
Representa un \dato\,\% adicional sobre el tiempo de la ISR. El resultado
\dato{} (se justifica / no se justifica) frente a la alternativa de que
cada driver implemente su propio \texttt{PORTB\_IRQHandler}, considerando
que a cambio cualquier módulo puede registrar un callback por pin sin tocar
el MCAL.

El período entre bits medido en M6 (\dato\,\us) corresponde a la velocidad
con la que se pasó la tarjeta en el ensayo; se retoma en el análisis de
latencia (M8).

\subsection{Asimetría entre las dos ISR del lector (M7 vs.\ M6)}

La ISR de ENABLE (M7) consume \dato\,\us\ en la apertura y \dato\,\us\ en
el cierre, es decir del orden de \dato\ veces el tiempo de la ISR de CLOCK
medida en M6. La diferencia se explica porque \texttt{IRQ\_CardEnable()}
reconfigura la interrupción de CLOCK mediante \texttt{gpioIRQ()}, operación
que implica lectura-modificación-escritura sobre el registro PCR del
puerto, limpieza del \textit{flag} en ISFR y habilitación en el NVIC.

Esta asimetría no representa un problema de diseño: la ISR larga ocurre
sólo dos veces por pasada, en momentos en que no hay bits llegando, mientras
que la ISR corta es la que se repite decenas de veces. El trabajo pesado
está ubicado donde no hay presión de tiempo.

\subsection{Validación de la política de prioridades (M8)}

Este es el resultado más relevante del informe. La Tabla~\ref{tab:latencia}
muestra una latencia máxima de \dato\,\us\ entre el flanco físico de CLOCK
y la entrada efectiva a \texttt{IRQ\_CardClk()}. Comparando ese valor contra
la duración de la ISR del SysTick medida en M1 (\dato\,\us):

\begin{itemize}
  \item Si la latencia máxima es notablemente menor a la duración de la ISR
        del SysTick, queda demostrado que PORTB efectivamente interrumpe al
        SysTick, y la asignación de prioridades (\texttt{APP\_NVIC\_PRIO\_}
        \texttt{SYSTICK = 1}) funciona como se diseñó.
  \item Si en cambio la latencia máxima se acerca a la duración de la ISR
        del SysTick, los flancos de CLOCK estarían esperando a que el
        SysTick termine, y el sistema estaría cerca de perder bits en
        pasadas rápidas.
\end{itemize}

El resultado obtenido indica que \dato[7cm], lo que \dato{} (valida /
invalida) experimentalmente la decisión de prioridades tomada en
\texttt{App\_Init()}.

Relacionando la latencia máxima con el período entre bits medido en M6
(\dato\,\us), el margen del sistema frente al peor caso es de
\dato\ veces: la tarjeta podría pasarse esa cantidad de veces más rápido
antes de que la latencia de interrupción comprometa la captura de un bit.

\subsection{Verificación del contrato de tiempo del encoder (M9)}

En reposo, el super-loop completa una vuelta cada \dato\,\us\ (Tabla~
\ref{tab:resultados}), muy por debajo del límite de 1 ms documentado en
\texttt{encoder.h}, así que el contrato de \texttt{getEncoderEvent()}
\dato{} (se cumple / no se cumple).

Durante la decodificación de una tarjeta el período sube a \dato\,\us,
porque el lazo ejecuta además \texttt{getCardData()} (búsqueda del
\textit{Start Sentinel} y verificación de paridad de cada carácter). Aun en
ese caso el período se mantiene \dato{} (por debajo / por encima) del
límite de 1 ms, lo que confirma que separar la captura (en la ISR) de la
decodificación (en el lazo principal) es una decisión de diseño correcta:
un retardo en la decodificación no causa pérdida de datos, porque los bits
ya están guardados en el arreglo.

\subsection{Conclusiones}

\begin{enumerate}
  \item El sistema \dato{} (cumple / no cumple) su restricción de tiempo
        real fundamental: la ISR del SysTick consume \dato\,\% de su
        período de 1 ms.
  \item El refresco del display es el componente dominante
        (\dato\,\% del tiempo de interrupción). Una futura optimización de
        rendimiento debería atacar el \textit{bit-banging} hacia los
        shift-registers, por ejemplo migrándolo a SPI por hardware.
  \item La latencia máxima medida para el lector (\dato\,\us) frente a la
        duración de la ISR del SysTick (\dato\,\us) \dato{} (confirma /
        refuta) que bajar la prioridad del SysTick fue la decisión
        correcta.
  \item Los mecanismos genéricos de despacho (SysTick y GPIO) suman
        \dato\,\us\ de sobrecarga total, un \dato\,\% del presupuesto de
        CPU. El equipo considera que este costo \dato{} (se justifica / no
        se justifica) frente a la ganancia en modularidad y portabilidad
        de la arquitectura en capas.
  \item Todas las mediciones corresponden a \texttt{-O0}. Se estima que con
        \texttt{-O2} los tiempos absolutos se reducirían \dato{}, aunque
        las proporciones relativas entre rutinas se mantendrían.
\end{enumerate}

\newpage
\appendix
% =============================================================================
\section{Capturas de osciloscopio}
% =============================================================================

\captura{m1_systick}{M1: pulso de la ISR completa del SysTick.}
\captura{m2_display_max}{M2: pulso del refresco del display, brillo máximo.}
\captura{m2_display_min}{M2: pulso del refresco del display, brillo mínimo.}
\captura{m3_encoder}{M3: pulso de la actualización del encoder.}
\captura{m4_timer}{M4: pulso del contador de milisegundos.}
\captura{m6_cardclk}{M6: ráfaga de pulsos de \texttt{IRQ\_CardClk} durante una pasada de tarjeta.}
\captura{m6_portdispatch}{M6: la misma ráfaga con el despachador de puerto incluido.}
\captura{m7_cardenable}{M7: pulsos de apertura y cierre de \texttt{IRQ\_CardEnable}.}
\captura{m8_latencia}{M8: flanco de CLOCK (CH1) contra entrada a la ISR (CH2).}
\captura{m9_looprepos}{M9: período del super-loop en reposo.}
\captura{m9_looptarjeta}{M9: período del super-loop decodificando una tarjeta.}

\end{document}
